\section{Discussion}
\label{sec:discussion}

\paragraph{Direction is the easy part.}
Seventeen surface cues decide \FW{} against \BW{} almost as reliably as a fine-tuned encoder: 8 confusions against 3, out of 380. What the cues cannot decide is the kind of relation. Does the added material narrow the meaning or paraphrase it? Does a swapped word break the relation or keep it? That is where most of the encoder's 0.21 gain in weighted F1 comes from (\cref{tab:slices}), and it is also where the encoder's own errors remain.

\paragraph{\SoftCons{} on its own is easy to game.}
A constant \EQ{} predictor scores a perfect 1.000. The cue model reaches 0.794 while being wrong in both directions on a quarter of the pairs. A system can be consistent by construction, through symmetric or sign-flipping inputs, without any grasp of the relation. We therefore read \SoftCons{} only next to \HardCons{} and macro F1, and we count the consistent-but-wrong pairs explicitly. We would suggest other participants do the same.

\paragraph{The encoder is consistent because it is correct.}
The trade between accuracy and consistency that we saw inside the classical family did not carry over. DeBERTa improved all three scores at once, and almost all of its consistent pairs are correct ones. For decoder models, which the pilot found far less coherent, this sharpens the question for Assignment 2: their consistency gap may be a correctness gap on the hard relations, and not a separate defect that needs its own fix.

\paragraph{How far the numbers travel.}
The development set is a little harder than held-out training folds, a fifth of its phrases also occur in training, and some gold labels look questionable. Differences of a few points between systems should be read with all of that in mind. The differences we draw conclusions from are 0.2 or more.
